% =====================================================================
%  c00.tex —— 标题页 / 摘要 / 译者说明（Lead 独占）
%  对应原文 PDF p.1（印刷 p.0）
% =====================================================================

\begin{titlepage}
\centering
\vspace*{3.2cm}

{\LARGE\bfseries 驹场镜像对称讲义\par}
\vspace{1.0cm}
{\Large\itshape Komaba Lectures on Mirror Symmetry\par}
\vspace{2.2cm}
{\Large 堀健太郎（Kentaro Hori）\par}
\vspace{0.6cm}
{\large 东京大学 Kavli 数物连携宇宙研究机构（Kavli IPMU, University of Tokyo, Japan）\par}
\vspace{3.2cm}
{\large 中译注本\par}
\vspace{0.4cm}
{\normalsize 由原文 LaTeX 排版重排为中文 \LaTeX{} 源码，附逐页校核与译者批注\par}
\vfill
{\small 底本：2016 年 8 月 26--28 日东京大学驹场校区（Komaba）讲义，共 49 页（含参考文献）。\par}
\vspace{0.4cm}
\end{titlepage}

\begin{center}
\begin{minipage}{0.92\textwidth}
\small
\noindent\textbf{摘要.}\quad 这是 2016 年 8 月 26--28 日在驹场（Komaba）所作讲座的讲义的一部分。

\medskip
\noindent\textit{Abstract.} A part of the notes for the lectures at Komaba on August 26--28, 2016.
\end{minipage}
\end{center}

\vspace{1.0cm}

\section*{译者说明}
\addcontentsline{toc}{section}{译者说明}

\begin{enumerate}[leftmargin=2em,itemsep=4pt]
\item \textbf{底本。}本译本依据 Kentaro Hori 的 \textit{Komaba Lectures on Mirror Symmetry}
  （2016 年 8 月东京大学驹场校区讲义；PDF 共 49 页，其中 p.1 为标题页、pp.2--48 为正文、
  p.49 为参考文献）。原讲义作者自述为“不完整版”（"A full version will be available soon"），
  故章节间详略不均，个别小标题下甚至暂无内容——这些地方本译本照原样保留并加批注，绝不补写。

\item \textbf{页码对应。}原文 PDF 的页脚印刷页码比 PDF 页号小 1，即
  \textbf{PDF 页号 = 印刷页码 + 1}。本译本用注释行 \verb|% [原文 p.N]| 标出每一个原文印刷页的起始位置，
  便于与底本对照。原文共 139 条编号公式，本译本公式编号与原文逐条一一对应
  （2.1--2.19、3.1--3.23、4.1--4.56、5.1--5.41），无增删、无重排。

\item \textbf{图。}原文 8 张矢量图（Figure 1--8）按 500 dpi 从底本裁取，图注由译者译出；
  图内的英文标注（如 "Phase I"）保持原样。原文无表格。

\item \textbf{术语。}专业术语以通行译法为准，首次出现时括注英文；模型名、缩写（GLSM、CFT、SCFT、
  T-duality、Calabi--Yau、Landau--Ginzburg、Fano、Bott--Morse 等）保留原文形式。
  详见工程目录中的 \verb|glossary.tsv|。

\item \textbf{批注。}译者的补充说明一律写成 LaTeX 注释 \verb|% [Tutor Note: ...]|，
  不进入排版；它们或指出原文笔误（正文照原样排，不改动），或补出作者跳过的推导，
  或标注易错点。全部批注汇总见随件《译注汇总.md》。

\item \textbf{参考文献。}按学术惯例不译，照原文转录（12 条），并保留原始出版信息。

\item \textbf{已知限制。}（一）原文为不完整讲义，第 5.5 节 "The $\CP^{N-1}$ Model" 小标题下无正文，
  系原文如此；（二）原文若干明显笔误（如 p.15 "empolyed"、p.29 "negtaive"、p.42 "ture" 等）
  按“正文不动、批注指出”的原则处理；（三）图示为位图裁取，放大后略见锯齿，但不影响内容。
\end{enumerate}

\clearpage
